% Options for packages loaded elsewhere
\PassOptionsToPackage{unicode}{hyperref}
\PassOptionsToPackage{hyphens}{url}
\documentclass[
  a4paper,
  10pt]{article}
\usepackage{xcolor}
\usepackage[margin=22mm]{geometry}
\usepackage{amsmath,amssymb}
\setcounter{secnumdepth}{-\maxdimen} % remove section numbering
\usepackage{iftex}
\ifPDFTeX
  \usepackage[T1]{fontenc}
  \usepackage[utf8]{inputenc}
  \usepackage{textcomp} % provide euro and other symbols
\else % if luatex or xetex
  \usepackage{unicode-math} % this also loads fontspec
  \defaultfontfeatures{Scale=MatchLowercase}
  \defaultfontfeatures[\rmfamily]{Ligatures=TeX,Scale=1}
\fi
\usepackage{lmodern}
\ifPDFTeX\else
  % xetex/luatex font selection
\fi
% Use upquote if available, for straight quotes in verbatim environments
\IfFileExists{upquote.sty}{\usepackage{upquote}}{}
\IfFileExists{microtype.sty}{% use microtype if available
  \usepackage[]{microtype}
  \UseMicrotypeSet[protrusion]{basicmath} % disable protrusion for tt fonts
}{}
\makeatletter
\@ifundefined{KOMAClassName}{% if non-KOMA class
  \IfFileExists{parskip.sty}{%
    \usepackage{parskip}
  }{% else
    \setlength{\parindent}{0pt}
    \setlength{\parskip}{6pt plus 2pt minus 1pt}}
}{% if KOMA class
  \KOMAoptions{parskip=half}}
\makeatother
\usepackage{longtable,booktabs,array}
\newcounter{none} % for unnumbered tables
\usepackage{calc} % for calculating minipage widths
% Correct order of tables after \paragraph or \subparagraph
\usepackage{etoolbox}
\makeatletter
\patchcmd\longtable{\par}{\if@noskipsec\mbox{}\fi\par}{}{}
\makeatother
% Allow footnotes in longtable head/foot
\IfFileExists{footnotehyper.sty}{\usepackage{footnotehyper}}{\usepackage{footnote}}
\makesavenoteenv{longtable}
\setlength{\emergencystretch}{3em} % prevent overfull lines
\providecommand{\tightlist}{%
  \setlength{\itemsep}{0pt}\setlength{\parskip}{0pt}}
\usepackage{bookmark}
\IfFileExists{xurl.sty}{\usepackage{xurl}}{} % add URL line breaks if available
\urlstyle{same}
\hypersetup{
  hidelinks,
  pdfcreator={LaTeX via pandoc}}

\author{}
\date{}

\begin{document}

\section{Complex Phase Structure Emerging from Nontrivial Quadratic
Closure}\label{complex-phase-structure-emerging-from-nontrivial-quadratic-closure}

\subsection{A Foundational Connection to Finite-Order Exchange Systems
and Discrete Born-Type
Weights}\label{a-foundational-connection-to-finite-order-exchange-systems-and-discrete-born-type-weights}

\textbf{Version:} English complete manuscript v1\\
\textbf{Date:} July 18, 2026\\
\textbf{Author:} Noriaki Kihara\\
\textbf{ORCID:} 0009-0004-6753-4020\\
\textbf{Version DOI:} 10.5281/zenodo.21422506\\
\textbf{Concept DOI:} 10.5281/zenodo.21422505\\
\textbf{Position in the series:} Additional interpretive paper deriving
from the quadratic-closure axiom in the Wave Information Readout series

\begin{center}\rule{0.5\linewidth}{0.5pt}\end{center}

\subsection{Abstract}\label{abstract}

This paper aims to provide a more foundational interpretation of the
previously reported result that finite-order closure of an
exchange-symmetric two-channel unitary operator yields the discrete
Born-type weights

\[
R_{n,m}=\cos^2\left(\frac{\pi m}{n}\right),
\qquad
1-R_{n,m}=\sin^2\left(\frac{\pi m}{n}\right).
\]

``Closure'' in this paper does not mean a thermodynamically closed
system isolated from the outside world. It means \textbf{algebraic
quadratic closure}: the conjugation-free algebraic quadratic form
\(Q_0(x)=\sum_kx_k^2\) closes nontrivially to zero.

The preceding paper treated exchange symmetry, unitarity, and the
finite-order condition within standard complex linear algebra. Here,
instead, we take as a minimal axiom over a commutative field \(K\)
containing the real field \(\mathbb R\) the nontrivial quadratic-closure
condition

\[
\boxed{
\sum_{k=1}^{N}x_k^2=0,
\qquad
(x_1,\ldots,x_N)\neq(0,\ldots,0),
\qquad
x_k\in K,
\qquad
\mathbb R\subseteq K
}.
\]

If every component were real, each term would be nonnegative and the
equation would have no nontrivial solution. The solution vector
\(\mathbf{x}:=(x_1,\ldots,x_N)\) must therefore lie outside
\(\mathbb R^N\); equivalently, at least one component is nonreal. In the
minimal two-component system

\[
x^2+y^2=0,
\qquad x\neq0,
\]

we have

\[
\left(\frac{y}{x}\right)^2=-1.
\]

Writing an element whose square is \(-1\) as \(i\) gives

\[
y=\pm ix.
\]

The imaginary unit is not introduced externally to describe phase; it is
required as the minimal extension direction that permits a nontrivial
sum of squares to close to zero.

We interpret this nonreal direction as a quarter turn. With continuity,
reversibility, and preservation of a positive-definite readout quantity
added as further conditions, it is extended to a general phase rotation
and connected to independently imposed exchange symmetry, the
finite-order condition

\[
\zeta^n=1,
\]

and the exchange-symmetric channel projections

\[
r=\frac{1+\zeta}{2},
\qquad
t=\frac{1-\zeta}{2}.
\]

For the finite-order root \(\zeta=e^{-2\pi im/n}\), taking squared
weights as a positive-definite readout recovers the previously reported
discrete Born-type sequence.

The rigorous derivation in this paper extends through the result that a
nontrivial zero sum of squares cannot be formed from real numbers alone
and that the minimal two-component system requires a
square-root-of-minus-one direction. The subsequent progression to
complex phase, an exchange-symmetric unitary representation, and
Born-type weights uses additional conditions: continuity, linearity,
reversibility, exchange symmetry, preservation of a positive-definite
readout, and squared readout. This paper is therefore not a derivation
of the complete Born rule. It classifies as derivation, additional
construction, and unresolved connection the relation between the nonreal
direction required by nontrivial quadratic closure and the finite-order
exchange theorem reported previously.

\textbf{Keywords:} quadratic closure, imaginary number, complex phase,
exchange symmetry, finite-order recurrence, Born rule, squared weight,
closed system, unitary operator

\begin{center}\rule{0.5\linewidth}{0.5pt}\end{center}

\section{1. Background and Objective}\label{1-background-and-objective}

\subsection{1.1 The Problem}\label{11-the-problem}

In standard quantum mechanics, states are represented as complex
vectors, time evolution is described by unitary operators, and
observation probabilities are given by squared amplitude magnitudes:

\[
P(A)=|\langle A|\psi\rangle|^2.
\]

For a two-state system,

\section{\$\$
\textbar\textbackslash psi\textbackslash rangle}\label{-psirangle}

\textbackslash cos\textbackslash phi,\textbar A\textbackslash rangle +
e\^{}\{i\textbackslash chi\}\textbackslash sin\textbackslash phi,\textbar B\textbackslash rangle,
\$\$

we have

\[
P(A)=\cos^2\phi,
\qquad
P(B)=\sin^2\phi.
\]

If complex numbers, conjugation, a positive-definite inner product,
unitarity, and squared magnitudes are all assumed together before the
\(\cos^2\) form is obtained, however, the origin of that square law may
not have been explained. Once complex amplitudes and conjugate products
have been adopted, two-state projections are already constrained
geometrically to have a \(\cos^2/\sin^2\) form.

The preceding paper, ``Emergence of Discrete Born-Type Weights in
Iterated Two-Channel Exchange Systems,'' did not take the Born rule as a
numerical search target. It discovered discrete Born-type weights
associated with finite-order recurrence in a model family addressing
wave-packet localization transfer, metastable two-state dynamics, weak
readout, and strong-observation selection. Its central result was that,
up to an overall phase, an exchange-symmetric two-channel unitary
operator can be written

\[
U=P_s+\zeta P_a,
\qquad |\zeta|=1,
\]

and that when the finite-order condition

\[
U^n=I
\]

selects

\[
\zeta=e^{-2\pi im/n},
\]

the channel weight is

\section{\$\$
\textbackslash left\textbar\textbackslash frac\{1+\textbackslash zeta\}\{2\}\textbackslash right\textbar\^{}2}\label{-leftfrac1zeta2right2}

\textbackslash cos\^{}2\textbackslash left(\textbackslash frac\{\textbackslash pi
m\}\{n\}\textbackslash right). \$\$

That result explains which phases and which squared weights are selected
by exact recurrence. A more anterior question nevertheless remains: why
should complex phase, conjugation, and unitarity be used at all?

\subsection{1.2 Question Addressed in This
Paper}\label{12-question-addressed-in-this-paper}

This paper asks:

\begin{quote}
Without taking complex numbers, conjugation, and unitarity as the first
axioms, how far can one construct the connection from a nontrivial
quadratic-closure condition to a square-root-of-minus-one direction,
phase rotation, exchange-symmetric finite-order structure, and discrete
Born-type weights?
\end{quote}

We distinguish four layers:

\[
\boxed{
\begin{gathered}
\text{nontrivial quadratic closure}
\Rightarrow\text{nonreal direction}
\quad[\text{derived}]\\
\text{nonreal direction}
+\{\text{continuity, reversibility, }Q_+\text{ preservation}\}\\
\longrightarrow\text{phase rotation and exchange representation}
\quad[\text{additional construction}]\\
\text{exchange symmetry}+\text{finite order}+\text{squared readout}\\
\Rightarrow\text{discrete Born-type weights}
\quad[\text{conditional derivation}]\\
Q_0=0\text{ and the dynamical connection to the exchange kernel}\\
[\text{connection problem C2}]
\end{gathered}
}.
\]

\subsection{1.3 Division of Roles Between the Two
Papers}\label{13-division-of-roles-between-the-two-papers}

This paper does not replace the preceding finite-order exchange theorem.

The role of the preceding paper is to derive discrete Born-type weights
rigorously within standard complex linear algebra from

\[
[U,X]=0,
\qquad
U^\dagger U=I,
\qquad
U^n=I.
\]

The role of the present paper is to show why a nonreal direction and
phase description are required by the preceding nontrivial closure
condition

\[
\sum_kx_k^2=0,
\]

and thereby to give a foundational interpretation of the mathematical
structure of the preceding paper.

The relation between the papers is thus

\[
\boxed{
\begin{gathered}
\text{this paper: derives the necessity of a nonreal direction}\\
\xrightarrow[\text{C2 remains unresolved}]{\text{additional construction}}\\
\text{preceding paper: derives Born-type weights from finite-order exchange}
\end{gathered}
}.
\]

The preceding paper is permanently available at Concept DOI
\url{https://doi.org/10.5281/zenodo.21422470}.

\begin{center}\rule{0.5\linewidth}{0.5pt}\end{center}

\section{2. Axiom of Nontrivial Quadratic
Closure}\label{2-axiom-of-nontrivial-quadratic-closure}

\subsection{2.1 Axiom}\label{21-axiom}

Let \(K\) be a commutative field containing the real field
\(\mathbb R\), and let there be finitely many components over it,

\[
x_k\in K,
\qquad
\mathbb R\subseteq K
\qquad(k=1,\ldots,N).
\]

We impose

\[
\boxed{
\sum_{k=1}^{N}x_k^2=0
}.
\]

To exclude the trivial solution alone, we also require

\[
\boxed{
(x_1,\ldots,x_N)\neq(0,\ldots,0)
}.
\]

We call this the \textbf{nontrivial quadratic-closure condition}.

Axiom A1 in the preceding work corresponds to the sector
\(K=\mathbb C\). Here we return the starting point to
\(\mathbb R\subseteq K\) in order to audit whether the conclusion arose
merely because the complex field had been fixed in advance. After
showing that two-component nontrivial closure generates a square root of
\(-1\) within \(K\), we recover the minimal subfield
\(\mathbb R(i)\cong\mathbb C\).

It is essential that the condition is not

\[
\sum_k|x_k|^2=0,
\]

but

\[
\sum_kx_k^2=0.
\]

The former is a positive-definite norm using complex conjugation and has
no nontrivial solution. The latter is a conjugation-free algebraic sum
of squares and may have nontrivial solutions when the number system is
extended beyond the reals.

\subsection{2.2 Triviality over the
Reals}\label{22-triviality-over-the-reals}

Suppose every component is real:

\[
x_k\in\mathbb R.
\]

Then

\[
x_k^2\ge0.
\]

For

\[
\sum_{k=1}^{N}x_k^2=0
\]

to hold, every term must therefore vanish:

\[
x_1=x_2=\cdots=x_N=0.
\]

\textbf{Proposition 2.1 (absence of nontrivial real solutions).}\\
Let \(N\ge1\). For \(x_k\in\mathbb R\), if

\[
\sum_{k=1}^{N}x_k^2=0,
\]

then

\[
x_k=0
\qquad(k=1,\ldots,N).
\]

Thus, if nontrivial quadratic closure is required, the solution vector
lies outside the real vector space \(\mathbb R^N\):

\[
\boxed{
\sum_kx_k^2=0,
\quad
(x_1,\ldots,x_N)\neq(0,\ldots,0)
\quad\Longrightarrow\quad
\mathbf{x}:=(x_1,\ldots,x_N)\in K^N\setminus\mathbb R^N
}.
\]

In components, this is equivalent to

\[
\boxed{
\exists j:\ x_j\notin\mathbb R
}.
\]

\subsection{2.3 Meaning of the
Proposition}\label{23-meaning-of-the-proposition}

The proposition does not require the entire extension field \(K\) to
equal the complex field. \(K\) may be an extension larger than
\(\mathbb C\).

The element \(i:=y/x\) obtained from two-component nontrivial closure
nevertheless satisfies \(i^2=-1\), so \(K\) contains the subfield

\[
\mathbb R(i).
\]

As a field over \(\mathbb R\), this is isomorphic to \(\mathbb C\). The
minimal commutative field extension obtained by adjoining a square root
of \(-1\) is therefore the complex field, while the solution vector of
nontrivial quadratic closure lies outside \(\mathbb R^N\).

In this paper, then, an imaginary number is not introduced arbitrarily
as a convenient notation for a wave function. It appears as

\begin{quote}
an extension direction required to close to zero a nontrivial sum of
squares that cannot close using positive real squares alone.
\end{quote}

\begin{center}\rule{0.5\linewidth}{0.5pt}\end{center}

\section{3. The Imaginary Direction in the Minimal Two-Component
System}\label{3-the-imaginary-direction-in-the-minimal-two-component-system}

\subsection{3.1 Two-Component Closure}\label{31-two-component-closure}

Consider the minimal nontrivial system

\[
x^2+y^2=0.
\]

If \(x\neq0\), then

\[
y^2=-x^2,
\]

and

\[
\left(\frac{y}{x}\right)^2=-1.
\]

Writing an element whose square is \(-1\) as \(i\),

\[
i^2=-1,
\]

we obtain

\[
\boxed{
y=\pm ix
}.
\]

\textbf{Proposition 3.1 (minimal two-component closure).}\\
In a commutative field extension satisfying \(x^2+y^2=0\) with
\(x\neq0\), the ratio \(y/x\) is a square root of \(-1\).

No complex conjugation occurs in this proposition.

If \(K\) contains \(\mathbb R\), the minimal subfield generated by
\(i:=y/x\) is

\[
\mathbb R(i)\cong\mathbb C.
\]

Thus, in two-component closure the complex direction appears not as an
externally supplied phase assumption, but as the minimal field extension
generated by a nontrivial solution.

\subsection{3.2 Interpretation as a Quarter
Turn}\label{32-interpretation-as-a-quarter-turn}

In the complex plane, multiplication of a quantity \(x\) along the real
axis by \(i\) corresponds to a 90-degree rotation:

\[
x\longmapsto ix.
\]

Therefore,

\[
y=\pm ix
\]

means that the two components are not merely positive or negative values
on the same line; they point in directions displaced by one quarter of a
cycle.

We call this relation a \textbf{quarter-period closure phase
difference}:

\[
\Delta\phi=\pm\frac{\pi}{2}.
\]

The phase difference was not assumed first to construct quadratic
closure. Rather, the nontriviality of quadratic closure requires a
square-root-of-minus-one direction, whose geometric representation is a
quarter turn.

\subsection{3.3 Sign and Opposite
Directions}\label{33-sign-and-opposite-directions}

The two solutions

\[
y=+ix,
\qquad
y=-ix
\]

have opposite directions of phase rotation:

\[
+\frac{\pi}{2},
\qquad
-\frac{\pi}{2}.
\]

This opposition may be related to the structure later represented by
complex conjugation.

We do not, however, claim to have derived the uniqueness of the
conjugation operation from this fact. What is rigorously obtained is
that a square root of \(-1\) has two signs, providing quarter-period
closure rotations in opposite directions.

\begin{center}\rule{0.5\linewidth}{0.5pt}\end{center}

\section{4. Nonreal Directions in a Multicomponent
System}\label{4-nonreal-directions-in-a-multicomponent-system}

\subsection{4.1 General Closure
Equation}\label{41-general-closure-equation}

Consider generally

\[
\sum_{k=1}^{N}x_k^2=0.
\]

Suppose at least one component \(x_j\) is nonzero. Since the square of a
real component is nonnegative, closure cannot be achieved using real
components alone.

For example, for \(a,b\in\mathbb R\), let

\[
x_1=a,
\qquad
x_2=b,
\qquad
x_3=i\sqrt{a^2+b^2}.
\]

Then

\[
a^2+b^2+
\left(i\sqrt{a^2+b^2}\right)^2
=0.
\]

This example shows that many registered real squares can be closed by
the square of a single imaginary direction.

\subsection{4.2 Necessary and Sufficient
Conditions}\label{42-necessary-and-sufficient-conditions}

The solution vector of nontrivial closure lies outside \(\mathbb R^N\).
In fixed closure coordinates, this is equivalent to at least one
component being nonreal.

The following stronger statements, however, do not hold in general:

\begin{itemize}
\tightlist
\item
  that there is exactly one nonreal component,
\item
  that a nonreal component is purely imaginary,
\item
  that each component separates into either a real or purely imaginary
  component, or
\item
  that the complex field is the only possible extension.
\end{itemize}

The weaker proposition sufficient for this paper is

\[
\boxed{
\text{every solution of nontrivial quadratic closure lies outside }\mathbb R^N
}.
\]

This motivates the introduction of nonreal degrees of freedom
representable as phase or rotation.

\subsection{4.3 Connection to
Unnamedness}\label{43-connection-to-unnamedness}

If no external name or absolute coordinate is assigned to each component
\(x_k\), and only the closure condition of the entire system is taken as
fundamental, then distinctions between components appear not as absolute
values but as relative allocations and phase relations of the registered
squares.

From this viewpoint,

\[
\sum_kx_k^2=0
\]

does not fix individual components as independent entities; it specifies
that all components constitute a single closed relational system.

\begin{center}\rule{0.5\linewidth}{0.5pt}\end{center}

\section{5. Distinguishing Quadratic Closure from Positive-Definite
Readout}\label{5-distinguishing-quadratic-closure-from-positive-definite-readout}

\subsection{5.1 Two Kinds of Quadratic
Form}\label{51-two-kinds-of-quadratic-form}

We clearly distinguish the following two forms.

\subsubsection{Quadratic form defining
closure}\label{quadratic-form-defining-closure}

\[
Q_0(x):=\sum_kx_k^2.
\]

It contains no conjugation.

The nontrivial closure condition is

\[
Q_0(x)=0,
\qquad (x_1,\ldots,x_N)\neq(0,\ldots,0).
\]

\subsubsection{Quadratic form for observation or
magnitude}\label{quadratic-form-for-observation-or-magnitude}

For complex components, define

\[
Q_+(x):=\sum_kx_k^*x_k
=\sum_k|x_k|^2.
\]

This form is positive definite:

\[
Q_+(x)\ge0.
\]

\subsection{5.2 Do Not Confuse the Closure Equation with the Norm
Equation}\label{52-do-not-confuse-the-closure-equation-with-the-norm-equation}

Nontrivial quadratic closure means

\[
Q_0(x)=0,
\]

not

\[
Q_+(x)=0.
\]

For a nontrivial state, one ordinarily has

\[
Q_0(x)=0,
\qquad
Q_+(x)>0.
\]

For the minimal two-component example

\[
(x,y)=(a,ia),
\]

we have

\[
Q_0=a^2+(ia)^2=0,
\]

but

\[
Q_+=|a|^2+|ia|^2=2|a|^2>0.
\]

This distinction means that a zero sum of squares need not be read as a
state in which all physical quantities have vanished. What vanishes is
the directional algebraic closure register; the positive-definite
readout may remain nonzero.

\subsection{5.3 Role of Conjugation}\label{53-role-of-conjugation}

In the construction presented here, conjugation is not a prerequisite
for the imaginary direction to exist.

The imaginary direction is already required by the nontriviality of

\[
x^2+y^2=0.
\]

Conjugation is introduced afterward as an operation that reverses the
nonreal direction,

\[
i\longmapsto-i,
\]

and converts

\[
z^*z
\]

into a real, positive-definite readout.

Separating the derived part from the added readout construction, the
logical order is therefore

\[
\boxed{
\text{nontrivial quadratic closure}
\Rightarrow
\text{necessity of a nonreal direction}
\qquad[\text{derived}]
}
\]

and

\[
\boxed{
\begin{gathered}
\text{nonreal direction}+\text{definition of the opposing reversal}\\
+\text{definition of positive-definite readout}
\Rightarrow Q_+\\
[\text{additional construction}]
\end{gathered}
}.
\]

This paper does not prove that the standard complex conjugation is
uniquely fixed by the nontrivial closure condition alone.

\begin{center}\rule{0.5\linewidth}{0.5pt}\end{center}

\section{6. Generalization to Phase
Rotation}\label{6-generalization-to-phase-rotation}

\subsection{6.1 Additional Construction from a Quarter-Period Closure
Phase to Continuous
Phase}\label{61-additional-construction-from-a-quarter-period-closure-phase-to-continuous-phase}

Minimal two-component closure requires the quarter-turn relation

\[
1,
\qquad
i.
\]

If continuity, associativity, reversibility, and preservation of a
positive-definite readout are added as conditions and this rotation is
extended to an iterable continuous group, one obtains phase rotations by

\[
e^{i\phi}.
\]

They satisfy

\[
e^{i\phi}e^{i\psi}=e^{i(\phi+\psi)}.
\]

The continuous phase group \(U(1)\) is not a consequence of the
nontrivial closure condition alone. We additionally require

\begin{enumerate}
\def\labelenumi{\arabic{enumi}.}
\tightlist
\item
  continuity of the rotation,
\item
  associativity of composition,
\item
  existence of an inverse, and
\item
  preservation of the positive-definite readout,
\end{enumerate}

and adopt the unit-circle phase as the resulting one-parameter
preservation group.

\subsection{6.2 Real Two-Dimensional
Representation}\label{62-real-two-dimensional-representation}

Without using complex numbers, the same transformation can be written as
the real two-dimensional rotation

\section{\$\$ R(\textbackslash phi)}\label{-rphi}

\textbackslash begin\{pmatrix\}
\textbackslash cos\textbackslash phi\&-\textbackslash sin\textbackslash phi\textbackslash{}
\textbackslash sin\textbackslash phi\&\textbackslash cos\textbackslash phi
\textbackslash end\{pmatrix\}. \$\$

Then

\[
R(\phi)^TR(\phi)=I,
\]

so the real two-dimensional norm is preserved.

Complex phase can therefore be regarded as a compressed representation
of real two-dimensional rotation.

The essential element in this paper is not the imaginary symbol itself,
but

\begin{quote}
extending the nonreal direction required by quadratic closure into a
reversible, norm-preserving rotational degree of freedom.
\end{quote}

\subsection{6.3 Finite-Order Phase}\label{63-finite-order-phase}

The condition that the rotation return to its initial state after
finitely many iterations is

\[
R(\phi)^n=I,
\]

or, in complex notation,

\[
\zeta^n=1.
\]

Thus,

\[
\zeta=e^{-2\pi im/n}.
\]

This integer-ratio phase gives a discrete closed orbit within the
continuous phase space.

\begin{center}\rule{0.5\linewidth}{0.5pt}\end{center}

\section{7. Connection to an Exchange-Symmetric Two-Channel
System}\label{7-connection-to-an-exchange-symmetric-two-channel-system}

\subsection{7.1 Two-Channel Exchange}\label{71-two-channel-exchange}

Let the exchange operator for the A/B channels be

\[
X=
\begin{pmatrix}
0&1\\
1&0
\end{pmatrix}.
\]

The symmetric and antisymmetric projectors are

\[
P_s=\frac{I+X}{2},
\qquad
P_a=\frac{I-X}{2}.
\]

Imposing

\[
[U,X]=0
\]

on an exchange-symmetric preserving map \(U\) prevents it from mixing
the symmetric and antisymmetric subspaces.

\subsection{7.2 Phase-Fixed
Representation}\label{72-phase-fixed-representation}

If a unitary representation is adopted as a linear reversible map
preserving the positive-definite readout, then, up to an overall phase,

\[
\boxed{
U=P_s+\zeta P_a,
\qquad |\zeta|=1
}
\]

can be written.

Here,

\begin{itemize}
\tightlist
\item
  the symmetric component is fixed, and
\item
  only the antisymmetric component undergoes phase rotation.
\end{itemize}

This is a \textbf{constructive hypothesis} placing the nonreal direction
required by nontrivial quadratic closure in the antisymmetric degree of
freedom of the two-channel exchange system.

\subsubsection{Connection Problem C2: Dynamical Nonpreservation of the
Quadratic-Closure
Surface}\label{connection-problem-c2-dynamical-nonpreservation-of-the-quadratic-closure-surface}

We now state explicitly the connection problem that remains between
nontrivial quadratic closure and exchange dynamics. Write the
two-channel state as

\[
x=(A,B)^T,
\qquad
Q_0(x)=x^Tx=A^2+B^2.
\]

For the exchange kernel

\[
U=P_s+\zeta P_a,
\]

the identities \(P_sP_a=0\), \(P_s^T=P_s\), and \(P_a^T=P_a\) give

\[
U^TU=P_s+\zeta^2P_a.
\]

Therefore,

\section{\$\$ \textbackslash boxed\{ Q\_0(Ux)}\label{-boxed-q_0ux}

x\^{}T(P\_s+\textbackslash zeta\^{}2P\_a)x \}. \$\$

For the nontrivial point on the closure surface

\[
x_0=(1,i)^T,
\qquad
Q_0(x_0)=1+i^2=0,
\]

one finds

\[
\boxed{
Q_0(Ux_0)=i(1-\zeta^2)
}.
\]

Thus, if \(\zeta^2\neq1\), \(Ux_0\) leaves the closure surface.
Moreover, this exchange kernel satisfies

\[
Q_0(Ux)=Q_0(x)
\]

for every state only when \(\zeta^2=1\). The previously reported
exchange kernel with general finite-order roots is therefore not a
self-map of the nontrivial quadratic-closure surface \(Q_0=0\).

This paper does not claim that the preceding exchange kernel directly
preserves \(Q_0=0\). We designate the construction of an exchange map
preserving the closure surface, a projection rule back onto that
surface, or an enlarged state space that preserves the closure condition
as \textbf{connection problem C2}. This is an open problem concerning a
dynamical connection between the first axiom and exchange dynamics; it
is independent of the validity of the finite-order exchange theorem
itself.

\subsection{7.3 Reconstruction in the Channel
Basis}\label{73-reconstruction-in-the-channel-basis}

Returning to the A/B basis gives

\section{\$\$ U}\label{-u}

\textbackslash frac12 \textbackslash begin\{pmatrix\}
1+\textbackslash zeta\&1-\textbackslash zeta\textbackslash{}
1-\textbackslash zeta\&1+\textbackslash zeta
\textbackslash end\{pmatrix\}, \$\$

and hence

\[
\boxed{
r=\frac{1+\zeta}{2},
\qquad
t=\frac{1-\zeta}{2}
}.
\]

This sum and difference are the interference between the fixed symmetric
phase \(1\) and the rotating antisymmetric phase \(\zeta\).

\begin{center}\rule{0.5\linewidth}{0.5pt}\end{center}

\section{8. Finite-Order Closure and Discrete Born-Type
Weights}\label{8-finite-order-closure-and-discrete-born-type-weights}

\subsection{8.1 Finite-Order Condition}\label{81-finite-order-condition}

After phase fixing, the condition for the operator to recur exactly
after \(n\) iterations is

\[
U^n=I.
\]

This is equivalent to

\[
\zeta^n=1.
\]

If \(n\) is the primitive order and \(\gcd(m,n)=1\), then

\[
\zeta=e^{-2\pi im/n}.
\]

\subsection{8.2 Squared Weights}\label{82-squared-weights}

For an A-basis input, the A/B channel amplitudes are

\[
r_{n,m}=\frac{1+e^{-2\pi im/n}}{2},
\]

\[
t_{n,m}=\frac{1-e^{-2\pi im/n}}{2}.
\]

Taking squared magnitudes as a positive-definite readout gives

\section{\$\$ \textbar r\_\{n,m\}\textbar\^{}2}\label{-r_nm2}

\section{\textbackslash frac\{1+\textbackslash cos(2\textbackslash pi
m/n)\}\{2\}}\label{frac1cos2pi-mn2}

\textbackslash cos\^{}2\textbackslash left(\textbackslash frac\{\textbackslash pi
m\}\{n\}\textbackslash right), \$\$

\section{\$\$ \textbar t\_\{n,m\}\textbar\^{}2}\label{-t_nm2}

\section{\textbackslash frac\{1-\textbackslash cos(2\textbackslash pi
m/n)\}\{2\}}\label{frac1-cos2pi-mn2}

\textbackslash sin\^{}2\textbackslash left(\textbackslash frac\{\textbackslash pi
m\}\{n\}\textbackslash right). \$\$

Therefore,

\section{\$\$ \textbackslash boxed\{ R\_\{n,m\}}\label{-boxed-r_nm}

\textbackslash cos\^{}2\textbackslash left(\textbackslash frac\{\textbackslash pi
m\}\{n\}\textbackslash right) \} \$\$

and

\section{\$\$ \textbackslash boxed\{ 1-R\_\{n,m\}}\label{-boxed-1-r_nm}

\textbackslash sin\^{}2\textbackslash left(\textbackslash frac\{\textbackslash pi
m\}\{n\}\textbackslash right) \} \$\$

are obtained.

\textbf{Theorem 8.1 (finite-order exchange weights under additional
conditions).}\\
Let \(U\) be a two-channel linear map satisfying the following
conditions.

\begin{enumerate}
\def\labelenumi{\arabic{enumi}.}
\tightlist
\item
  \(U\) is an invertible map preserving the positive-definite readout
  \(Q_+\).
\item
  \(U\) commutes with the exchange operator \(X\).
\item
  The overall phase is fixed so that the eigenvalue of the symmetric
  sector is \(1\).
\item
  \(U^n=I\), and the primitive order of the antisymmetric sector is
  \(n\).
\item
  The A/B channel weights are defined by a positive-definite squared
  readout of the amplitudes.
\end{enumerate}

Then, for coprime integers \(m,n\),

\[
U=P_s+e^{-2\pi im/n}P_a,
\]

and the A/B channel weights for an A-basis input are

\[
\boxed{
R_{n,m}=\cos^2\left(\frac{\pi m}{n}\right),
\qquad
1-R_{n,m}=\sin^2\left(\frac{\pi m}{n}\right)
}.
\]

\emph{Proof.} Conditions 1--3 give \(U=P_s+\zeta P_a\) with
\(|\zeta|=1\). Condition 4 gives \(\zeta^n=1\), and hence
\(\zeta=e^{-2\pi im/n}\). Returning to the A/B basis gives the
amplitudes \((1\pm\zeta)/2\), and the positive-definite squared readout
in condition 5 gives the displayed result. \(\square\)

This theorem restates the finite-order exchange theorem of the preceding
paper so that the present paper can be read independently. Nontrivial
quadratic closure does not derive conditions 1--5 themselves; it
supplies the nonreal direction required for the complex-phase
representation. The dynamical connection between the nontrivial closure
surface and \(U\) is C2 and is not included among the assumptions of
this theorem.

\subsection{8.3 Foundational
Interpretation}\label{83-foundational-interpretation}

The preceding paper obtained the formulas above as a finite-order
theorem for an exchange-symmetric unitary kernel.

Here we interpret their more anterior stage through nontrivial closure,

\[
\boxed{
\sum_kx_k^2=0
}.
\]

The logical sequence is therefore not a single chain of unconditional
implications, but is classified into four layers:

\[
\boxed{
\begin{gathered}
\text{nontrivial quadratic closure}
\Rightarrow\text{a square-root-of-minus-one direction}
\quad[\text{derived}]\\
\text{nonreal direction}
+\{\text{continuity, reversibility, }Q_+\text{ preservation}\}\\
\longrightarrow U(1)\text{ phase rotation}
\quad[\text{additional construction}]\\
U(1)\text{ phase}+\text{exchange symmetry}+\text{finite order}\\
+\text{squared readout}
\Rightarrow\text{discrete Born-type weights}
\quad[\text{conditional derivation}]\\
Q_0=0\text{ and the dynamical connection to the exchange kernel }U\\
[\text{connection problem C2}]
\end{gathered}
}.
\]

\begin{center}\rule{0.5\linewidth}{0.5pt}\end{center}

\section{9. Relation to the Born
Rule}\label{9-relation-to-the-born-rule}

\subsection{9.1 What Has Been
Explained}\label{91-what-has-been-explained}

Taken together, this paper and the preceding paper explain the
following.

\begin{enumerate}
\def\labelenumi{\arabic{enumi}.}
\tightlist
\item
  A nontrivial zero sum of squares cannot be formed from real numbers
  alone. {[}Derived{]}
\item
  Minimal two-component closure requires a direction whose square is
  \(-1\). {[}Derived{]}
\item
  This direction can be represented as a quarter turn. {[}Standard
  geometric interpretation{]}
\item
  Continuity, reversibility, and \(Q_+\) preservation extend it to phase
  rotation. {[}Additional construction{]}
\item
  The phase is placed in the antisymmetric sector of an
  exchange-symmetric two-channel system. {[}Constructive hypothesis{]}
\item
  The finite-order condition discretizes the phase to an integer ratio.
  {[}Conditional derivation{]}
\item
  The squared weights of the interference amplitudes returned to the A/B
  basis are \(\cos^2/\sin^2\). {[}Conditional derivation{]}
\end{enumerate}

In this sense, the present paper offers

\[
\text{nontrivial quadratic closure}
\]

as a candidate root preceding the Born-type weights.

\subsection{9.2 What Has Not Yet Been
Explained}\label{92-what-has-not-yet-been-explained}

This paper does not prove

\begin{enumerate}
\def\labelenumi{\arabic{enumi}.}
\tightlist
\item
  the complete Born rule of standard quantum mechanics,
\item
  why a positive-definite readout can be identified with an experimental
  frequency,
\item
  how single-trial selection into A or B occurs,
\item
  convergence of repeated-trial frequencies to \(|r|^2\) and \(|t|^2\),
\item
  that the complex field is the unique possible extension,
\item
  that standard complex conjugation is uniquely derived from nontrivial
  quadratic closure alone,
\item
  that the unitary group is uniquely derived without additional
  conditions,
\item
  projection measures in arbitrary dimension, or
\item
  that the preceding exchange kernel preserves the nontrivial
  quadratic-closure surface \(Q_0=0\).
\end{enumerate}

The conclusion is therefore not

\begin{quote}
The Born rule has been derived completely.
\end{quote}

More precisely, it is

\begin{quote}
Nontrivial quadratic closure necessitates complex-phase structure, and
when connected to an exchange-symmetric finite-order system, it yields
discrete Born-type squared weights.
\end{quote}

\subsection{9.3 Significance for
Circularity}\label{93-significance-for-circularity}

A conventional circular explanation has the form

\[
\text{complex amplitude}
\rightarrow
\text{conjugate product}
\rightarrow
\cos^2.
\]

The construction here is

\[
\begin{aligned}
\text{nontrivial quadratic closure}
&\rightarrow\text{imaginary direction}
\rightarrow\text{phase rotation}\\
&\rightarrow\text{finite closure}
\rightarrow\text{channel amplitude}
\rightarrow\text{squared readout}.
\end{aligned}
\]

The remaining inputs in this construction are

\begin{itemize}
\tightlist
\item
  the positive-definite readout \(Q_+\),
\item
  adoption of that readout as the channel weight, and
\item
  an additional physical process connecting that weight to experimental
  frequency.
\end{itemize}

This paper therefore does not derive the origin of the squared norm
itself or the Born frequency law.

What it derives, or determines conditionally, is

\begin{itemize}
\tightlist
\item
  that nontrivial quadratic closure requires a nonreal direction,
\item
  that the finite-order condition selects integer-ratio phases, and
\item
  that projection of those phases onto an exchange-symmetric A/B basis
  gives discrete \(\cos^2/\sin^2\) weights.
\end{itemize}

Thus, the inputs removed are the external provision of a nonreal
direction and the external selection of the phase angle that closes. The
inputs remaining are positive-definite squared readout and its frequency
interpretation. This distinction precisely delimits the noncircular part
of the derivation.

\begin{center}\rule{0.5\linewidth}{0.5pt}\end{center}

\section{10. Meaning of the Closed-System
Assumption}\label{10-meaning-of-the-closed-system-assumption}

\subsection{10.1 Definition of
``Closure''}\label{101-definition-of-closure}

Closure in this paper does not immediately mean a physical container
completely isolated in space from the outside world.

Its minimal meaning is that the relation

\[
\sum_kx_k^2=0
\]

is complete within the set of components under consideration.

Thus, closure means

\begin{quote}
algebraic closure in which no undefined component outside the system is
needed to evaluate the sum of the registered quantities under
consideration.
\end{quote}

\subsection{10.2 Distinction from Physical
Closure}\label{102-distinction-from-physical-closure}

Algebraic closure is different from thermodynamic, causal, or
experimental closure.

This paper does not claim that a real quantum experiment is completely
isolated from the outside world.

It considers instead the possibility that an enlarged effective system
containing, over finite time and finite degrees of freedom,

\begin{itemize}
\tightlist
\item
  the target channels,
\item
  the measurement apparatus, and
\item
  the environment
\end{itemize}

may admit an internal closure relation.

\subsection{10.3 Why This Is a Separate
Paper}\label{103-why-this-is-a-separate-paper}

The finite-order exchange theorem in the preceding paper holds using
standard unitary-operator theory alone.

Nontrivial quadratic closure in this paper, by contrast, contains
foundational assumptions about the number system and the meaning of
physical closure.

Combining the two in one paper would conflate the validity of the
finite-order theorem with the physical validity of the closure axiom.

This paper therefore leaves the preceding work unchanged and presents
the foundational argument independently as an additional interpretation.

\begin{center}\rule{0.5\linewidth}{0.5pt}\end{center}

\section{11. Falsifiable Tests}\label{11-falsifiable-tests}

\subsection{11.1 Reconstruction in a Real
Representation}\label{111-reconstruction-in-a-real-representation}

Reconstruct the route from

\[
\sum_kx_k^2=0
\]

to the finite-order weights of the exchange system using only real
two-dimensional rotations and no complex-number notation.

This would confirm that the result does not depend on symbolic
manipulations specific to complex notation.

\subsection{11.2 Comparison of Preserved
Forms}\label{112-comparison-of-preserved-forms}

Compare the following three conditions:

\begin{enumerate}
\def\labelenumi{\arabic{enumi}.}
\tightlist
\item
  the conjugation-free quadratic closure \(Q_0(x)=\sum x_k^2\),
\item
  the positive-definite norm \(Q_+(x)=\sum|x_k|^2\), and
\item
  the indefinite-metric form \(Q_G(x)=x^TGx\).
\end{enumerate}

Determine which form generates the same finite-order sequence as the
preceding exchange kernel under the fewest assumptions.

\subsection{11.3 Uniqueness of the Conjugation
Map}\label{113-uniqueness-of-the-conjugation-map}

For a map that reverses a nonreal direction and gives a
positive-definite readout, examine whether standard complex conjugation
is determined uniquely when one imposes

\begin{itemize}
\tightlist
\item
  additivity,
\item
  multiplicative order,
\item
  involutivity, and
\item
  fixation of real numbers.
\end{itemize}

\subsection{11.4 Derivation of the Preservation
Group}\label{114-derivation-of-the-preservation-group}

Classify the linear transformation group preserving the nontrivial
closure set

\section{\$\$ \textbackslash mathcal C}\label{-mathcal-c}

\textbackslash left\{x\textbackslash neq0\textbackslash mid\textbackslash sum\_kx\_k\^{}2=0\textbackslash right\}.
\$\$

Then clarify under which conditions the subgroup that also preserves an
added positive-definite readout becomes an orthogonal or unitary group.

\subsection{11.5 Connection to the Frequency
Law}\label{115-connection-to-the-frequency-law}

For the finite-order exchange weights

\section{\$\$ R\_\{n,m\}}\label{-r_nm}

\textbackslash cos\^{}2\textbackslash left(\textbackslash frac\{\textbackslash pi
m\}\{n\}\textbackslash right), \$\$

perform repeated trials with observational backreaction and phase
fluctuations, and determine the conditions under which

\[
f_A\to R_{n,m},
\qquad
f_B\to1-R_{n,m}.
\]

\begin{center}\rule{0.5\linewidth}{0.5pt}\end{center}

\section{12. Discussion}\label{12-discussion}

\subsection{12.1 Imaginary Numbers as Directions Required for
Closure}\label{121-imaginary-numbers-as-directions-required-for-closure}

The simplest consequence of this paper is that requiring a nontrivial
solution of

\[
\sum_kx_k^2=0
\]

is impossible using real numbers alone.

For two components, one necessarily obtains

\[
y=\pm ix.
\]

An imaginary number can therefore be interpreted not as an arbitrary
symbol appended to the real numbers, but as the minimal opposing
direction required to realize quadratic closure nontrivially.

\subsection{12.2 Conjugation Belongs to Readout, Not to
Closure}\label{122-conjugation-belongs-to-readout-not-to-closure}

The nontrivial closure equation requires no conjugation.

The square-root-of-minus-one direction is already required by

\[
x^2+y^2=0.
\]

Conjugation is needed later, when that nonreal direction is reversed to
read the positive-definite quantity

\[
z^*z.
\]

Thus,

\begin{itemize}
\tightlist
\item
  the necessity of imaginary numbers,
\item
  the role of conjugation, and
\item
  the Born frequency law
\end{itemize}

are distinct problems.

\subsection{12.3 Roles of the Two Kinds of
Closure}\label{123-roles-of-the-two-kinds-of-closure}

Two kinds of closure appear in this paper, but they do not have the same
role.

The first is the algebraic quadratic closure

\[
\sum_kx_k^2=0.
\]

The second is the finite-order phase closure

\[
\zeta^n=1.
\]

The first closure rigorously requires a nonreal direction. The second
selects integer-ratio phases from the phase rotations supplied by the
additional construction.

Returning that phase to the exchange-symmetric A/B basis gives the
interference amplitudes

\[
\frac{1\pm\zeta}{2},
\]

whose positive-definite squared readouts are

\[
\cos^2\left(\frac{\pi m}{n}\right),
\qquad
\sin^2\left(\frac{\pi m}{n}\right).
\]

The condition that directly derives the discrete Born-type weights is
therefore

\[
\boxed{
\text{phase closure}
+
\text{exchange projection}
+
\text{positive-definite squared readout}
\Rightarrow
\text{discrete Born-type weights}
}.
\]

Nontrivial quadratic closure provides a foundation for the nonreal
direction of the complex-phase structure entering this derivation. It
does not automatically follow, however, that the second exchange
dynamics preserves the first closure surface; that dynamical connection
remains problem C2.

\subsection{12.4 Scope of Derivation and Additional
Conditions}\label{124-scope-of-derivation-and-additional-conditions}

The central result derived directly here is that nontrivial quadratic
closure requires a direction beyond the real number system, even if
complex numbers are not postulated initially as an axiom of quantum
theory.

We do not claim to have uniquely derived all the mathematics of standard
quantum mechanics from nontrivial closure.

In particular, further conditions are still required between

\begin{itemize}
\tightlist
\item
  the complex field,
\item
  standard conjugation,
\item
  a positive-definite inner product,
\item
  the unitary group, and
\item
  the Born frequency law.
\end{itemize}

This paper does not conceal those conditions; it separates the derived
part from the interpretive connection.

\begin{center}\rule{0.5\linewidth}{0.5pt}\end{center}

\section{13. Conclusion}\label{13-conclusion}

This paper began from the nontrivial quadratic-closure condition

\[
\boxed{
\sum_{k=1}^{N}x_k^2=0,
\qquad
(x_1,\ldots,x_N)\neq(0,\ldots,0),
\qquad
x_k\in K,
\quad
\mathbb R\subseteq K
}.
\]

If every component is real, each squared term is nonnegative, so the
equation has only the trivial solution. The solution vector of
nontrivial closure therefore lies outside \(\mathbb R^N\). In fixed
closure coordinates this is equivalent to at least one component being
nonreal.

In the minimal two-component system,

\[
x^2+y^2=0,
\qquad x\neq0,
\]

we obtain

\[
\boxed{
y=\pm ix
}.
\]

No complex conjugation is used in this result. A
square-root-of-minus-one direction is required to form a nontrivial zero
sum of squares.

If continuity, reversibility, and preservation of a positive-definite
readout are added to this nonreal direction to extend it to a phase
rotation, and the resulting phase is placed in the antisymmetric mode of
an exchange-symmetric two-channel system, one obtains

\[
U=P_s+\zeta P_a.
\]

The finite-order condition

\[
\zeta^n=1
\]

selects the discrete phase

\[
\zeta=e^{-2\pi im/n}.
\]

The interference amplitudes in the A/B basis are

\[
r=\frac{1+\zeta}{2},
\qquad
t=\frac{1-\zeta}{2},
\]

and their positive-definite squared readouts are

\section{\$\$ \textbackslash boxed\{
\textbar r\textbar\^{}2}\label{-boxed-r2}

\textbackslash cos\^{}2\textbackslash left(\textbackslash frac\{\textbackslash pi
m\}\{n\}\textbackslash right) \} \$\$

and

\section{\$\$ \textbackslash boxed\{
\textbar t\textbar\^{}2}\label{-boxed-t2}

\textbackslash sin\^{}2\textbackslash left(\textbackslash frac\{\textbackslash pi
m\}\{n\}\textbackslash right) \}. \$\$

Separating derivation, additional construction, and the unresolved
connection, the relation established here is

\[
\boxed{
\begin{gathered}
\text{nontrivial quadratic closure}
\Rightarrow\text{nonreal direction}
\quad[\text{derived}]\\
\text{nonreal direction}
+\{\text{continuity, reversibility, }Q_+\text{ preservation}\}\\
\longrightarrow\text{phase rotation}
\quad[\text{additional construction}]\\
\text{phase rotation}+\text{exchange symmetry}+\text{finite order}\\
+\text{squared readout}
\Rightarrow\text{discrete Born-type weights}
\quad[\text{conditional derivation}]\\
Q_0=0\text{ and its dynamical connection to the exchange kernel}\\
[\text{connection problem C2}]
\end{gathered}
}.
\]

This is not a derivation of the complete Born rule. What has been
derived rigorously is that nontrivial quadratic closure cannot be
achieved using real numbers alone, that the minimal two-component system
requires a square-root-of-minus-one direction, and that combination with
the preceding finite-order exchange theorem recovers discrete Born-type
weights.

Future problems include unique derivations of conjugation,
positive-definite readout, unitarity, and a frequency law, as well as
construction of an exchange map that preserves \(Q_0=0\) or a projection
rule onto the closure surface. The central significance of this paper is
that complex phase is positioned not as a tool specific to quantum
theory assumed from the beginning, but as a necessary condition for
nontrivial quadratic closure.

\begin{center}\rule{0.5\linewidth}{0.5pt}\end{center}

\section{References}\label{references}

\subsection{Self-Citations}\label{self-citations}

\begin{enumerate}
\def\labelenumi{\arabic{enumi}.}
\tightlist
\item
  Noriaki Kihara, ``Discovery of Finite-Order Resonances in Iterated
  Exchange Scattering: Identification of the Origin of Peaks Near the
  Fine-Structure Values 137 and 128 and a Reproducible Wave-Packet
  Mathematical Model,'' Version DOI: 10.5281/zenodo.21421367, Concept
  DOI: 10.5281/zenodo.21421366, 2026.
\item
  Noriaki Kihara, ``First-Principles Derivation of the Minimal
  Nontrivial 137 Cells in a Closed Two-Wave Exchange System: Two
  Principal Candidates, a Conditional Dual Attractor, and a
  Correspondence Hypothesis with the Fine-Structure Constant,'' v3.0,
  2026. Repository version:
  \href{../20260717/20260717_complete_paper_v3.md}{\texttt{20260717\_complete\_paper\_v3.md}}.
  {[}Japanese{]}.
\item
  Noriaki Kihara, ``Emergence of Discrete Born-Type Weights in Iterated
  Two-Channel Exchange Systems: A Finite-Order Recurrence Law from
  Wave-Packet Localization Transfer, Metastable Two-State Dynamics, and
  Observation Selection,'' Version DOI: 10.5281/zenodo.21422471, Concept
  DOI: 10.5281/zenodo.21422470, 2026.
\item
  Noriaki Kihara, ``Experimental Specification v1 for Low-Localization
  and Harmonic Transfer Readout in Fermion-Like Collisions Using an
  Exchange-Interference Scattering Matrix,'' 2026. {[}Japanese{]}.
\item
  Noriaki Kihara, ``Preliminary Experimental Summary v1 of the
  Acceleration Basis and Localization Exchange in Fermion-Like
  Collisions Using an Exchange-Interference Scattering Matrix,'' Zenodo
  Concept DOI: 10.5281/zenodo.21333766, 2026. {[}Japanese{]}.
\item
  Noriaki Kihara, ``Preliminary Experimental Summary v1 of C Weak
  Readout and D Strong-Observation Selection at a White-Cat, Black-Cat,
  and Gray-Cat Metastable Interface,'' Zenodo Concept DOI:
  10.5281/zenodo.21353208, 2026. {[}Japanese{]}.
\end{enumerate}

\subsection{External References}\label{external-references}

\begin{enumerate}
\def\labelenumi{\arabic{enumi}.}
\setcounter{enumi}{6}
\tightlist
\item
  Max Born, ``Zur Quantenmechanik der Stoßvorgänge,'' \emph{Zeitschrift
  für Physik} \textbf{37}, 863--867 (1926). DOI: 10.1007/BF01397477.
\item
  Andrew M. Gleason, ``Measures on the Closed Subspaces of a Hilbert
  Space,'' \emph{Journal of Mathematics and Mechanics} \textbf{6},
  885--893 (1957). DOI: 10.1512/iumj.1957.6.56050.
\end{enumerate}

\begin{center}\rule{0.5\linewidth}{0.5pt}\end{center}

\section{Appendix A. Classification of
Claims}\label{appendix-a-classification-of-claims}

{\def\LTcaptype{none} % do not increment counter
\small
\begin{longtable}[]{@{}p{0.54\linewidth}p{0.39\linewidth}@{}}
\toprule\noalign{}
Claim & Classification \\
\midrule\noalign{}
\endhead
\bottomrule\noalign{}
\endlastfoot
Requiring a nontrivial solution of \(\sum_kx_k^2=0\) & First axiom of
this paper \\
No nontrivial solution of \(\sum_kx_k^2=0\) exists using only real
components & Derived proposition \\
A nontrivial solution vector belongs to \(K^N\setminus\mathbb R^N\) &
Derived proposition \\
In fixed closure coordinates, at least one component is nonreal &
Equivalent component representation of the preceding proposition \\
In the two-component system \(x^2+y^2=0\), \(y/x\) is a square root of
\(-1\) & Derived proposition \\
The minimal subfield generated by two-component nontrivial closure is
\(\mathbb R(i)\cong\mathbb C\) & Derived consequence \\
If a square root of \(-1\) is denoted by \(i\), then \(y=\pm ix\) &
Derived consequence \\
Reading \(\pm i\) as a quarter turn & Standard geometric
interpretation \\
Extending the nonreal direction to the continuous phase group \(U(1)\) &
Construction adding continuity, reversibility, and preservation \\
Standard complex conjugation is uniquely derived from nontrivial
quadratic closure & Not derived \\
The unitary group is uniquely derived from nontrivial quadratic closure
& Not derived \\
The exchange-symmetric finite-order kernel gives discrete Born-type
weights & Derived in the preceding paper \\
The preceding exchange kernel preserves the nontrivial quadratic-closure
surface \(Q_0=0\) & False in general; connection problem C2 \\
Nontrivial quadratic closure is the physical origin of complex-phase
structure & Foundational interpretive hypothesis \\
Discrete Born-type weights give experimental frequencies & Untested \\
The complete Born rule has been derived & Not derived \\
\end{longtable}
}

\begin{center}\rule{0.5\linewidth}{0.5pt}\end{center}

\section{Appendix B. Minimal Derivation
Sequence}\label{appendix-b-minimal-derivation-sequence}

\subsection{B.1 Nontrivial Closure}\label{b1-nontrivial-closure}

\[
\sum_{k=1}^{N}x_k^2=0,
\qquad
(x_1,\ldots,x_N)\neq(0,\ldots,0),
\qquad
x_k\in K,
\qquad
\mathbb R\subseteq K.
\]

\subsection{B.2 Impossibility Using Real Numbers
Alone}\label{b2-impossibility-using-real-numbers-alone}

\[
x_k\in\mathbb R
\quad\Longrightarrow\quad
x_k^2\ge0.
\]

Therefore,

\[
\sum_kx_k^2=0
\quad\Longrightarrow\quad
x_k=0\ \forall k.
\]

For a nontrivial solution, first in vector notation,

\[
\mathbf{x}\in K^N\setminus\mathbb R^N.
\]

The equivalent component representation in fixed closure coordinates is

\[
\exists j:\ x_j\notin\mathbb R.
\]

\subsection{B.3 Minimal Two-Component
System}\label{b3-minimal-two-component-system}

From

\[
x^2+y^2=0,
\qquad x\neq0,
\]

we have

\[
\left(\frac{y}{x}\right)^2=-1.
\]

If a square root of \(-1\) is denoted by \(i\),

\[
y=\pm ix.
\]

\subsection{B.4 Additional Construction and Phase
Closure}\label{b4-additional-construction-and-phase-closure}

The following is not a direct consequence of nontrivial quadratic
closure alone. We add continuity, reversibility, preservation of a
positive-definite readout, and exchange symmetry, representing the
nonreal direction as the antisymmetric-sector phase \(\zeta\). Imposing
finite order on this additional construction gives

\[
\zeta^n=1
\quad\Longrightarrow\quad
\zeta=e^{-2\pi im/n}.
\]

\subsection{B.5 Exchange Projection}\label{b5-exchange-projection}

\section{\$\$ U=P\_s+\textbackslash zeta P\_a}\label{-up_szeta-p_a}

\textbackslash frac12 \textbackslash begin\{pmatrix\}
1+\textbackslash zeta\&1-\textbackslash zeta\textbackslash{}
1-\textbackslash zeta\&1+\textbackslash zeta
\textbackslash end\{pmatrix\}. \$\$

\subsection{B.6 Discrete Born-Type
Weights}\label{b6-discrete-born-type-weights}

\[
r=\frac{1+\zeta}{2},
\qquad
t=\frac{1-\zeta}{2},
\]

and therefore

\section{\$\$ \textbackslash boxed\{
\textbar r\textbar\^{}2}\label{-boxed-r2-1}

\textbackslash cos\^{}2\textbackslash left(\textbackslash frac\{\textbackslash pi
m\}\{n\}\textbackslash right) \} \$\$

\section{\$\$ \textbackslash boxed\{
\textbar t\textbar\^{}2}\label{-boxed-t2-1}

\textbackslash sin\^{}2\textbackslash left(\textbackslash frac\{\textbackslash pi
m\}\{n\}\textbackslash right) \}. \$\$

\end{document}
